% !Mode:: "TeX:UTF-8"

\makeatletter % 允许在正文中使用带有 @ 的内部变量

%%% 此部分需要自行填写: 中文摘要及关键词 

%%%---- 1. 诚信承诺书 ------------------------------------%
\newpage
\thispagestyle{empty}
\bgroup
  \setlength{\parindent}{0em}
  \linespread{1.5}\selectfont

  \vspace*{1em}
  {\centering \songti\zihao{3} 北方民族大学本科毕业论文（设计）诚信承诺书 \par}

  \vspace{1.5em}

  \noindent
  \begin{tabular}{|p{0.18\textwidth}|p{0.28\textwidth}|p{0.15\textwidth}|p{0.28\textwidth}|}
  \hline
  \centering 学生姓名 & \centering \@author & \centering 年级 & \centering\arraybackslash 2022 级 \\
  \hline
  \centering 所学专业 & \centering \the\Cmajor & \centering 学号 & \centering\arraybackslash \the\CstudentNum \\
  \hline
  \centering 所在学院 & \multicolumn{3}{c|}{\the\Ccoursename} \\
  \hline
  \multicolumn{4}{|p{\dimexpr\textwidth-2\tabcolsep-2\arrayrulewidth\relax}|}{
    \def\arraystretch{1.5}
    
    \vspace{1em}
    {\centering \heiti\zihao{3} 学生承诺 \par}
    
    \vspace{1.5em}
    {
      \songti\zihao{4} 
      \setlength{\parindent}{2em}
      \noindent \textbf{本人郑重承诺和声明：}\par
      \indent 我承诺在毕业论文（设计）过程中严格遵守学校有关规定，在指导教师的安排与指导下独立完成所规定的毕业论文（设计）工作，决不弄虚作假，不请别人代做毕业论文（设计）或抄袭别人的成果。所撰写的毕业论文或毕业设计是在指导老师的指导下自主完成，文中所有引文或引用数据、图表均注解并说明来源，本人愿意为由此引起的后果承担责任。\par
    }
    
    \vspace{5em}
    \begin{flushright}
      {\songti\zihao{4} 
      学生（签名）：\makebox[8em]{\hrulefill} \hspace{2em} \par
      \vspace{1.5em}
      2026 年 05 月 \hspace{1em} 日 \hspace{2em} \par
      }
    \end{flushright}
    \vspace{1.5em}
  } \\
  \hline
  \end{tabular}

  \vfill

  \begin{flushright}
    {\zihao{4} 教务处制}
  \end{flushright}
\egroup
\clearpage

%%======摘要===========================%
\begin{cnabstract}
\thispagestyle{empty}

在现代实时计算机图形学中，传统的纯光栅化渲染在处理全局光照、精确物理软阴影与复杂镜面反射时存在固有的物理局限性，而离线级别的纯路径追踪又难以满足交互式应用的实时帧率需求。针对这一“画质与算力”的矛盾，本文基于现代显式图形 API Vulkan 1.3，设计并实现了一套工业级的高性能实时光线追踪混合渲染引擎。

本文的核心工作包括三个方面。首先，针对 Vulkan 繁琐的显存状态同步痛点，设计了数据驱动的 Render Graph 调度架构，通过有向无环图的拓扑分析实现了管线屏障的自动推导与显存复用优化。其次，构建了高效的混合光照管线，将几何光栅化与硬件光线追踪深度解耦。在延迟渲染 G-Buffer 的基础上，结合偏导数面法线偏移与 Ray Query 扩展实现了精准无伪影的软硬阴影，并利用 RT Pipeline 全面求解 Cook-Torrance 微面元模型，完成了基于物理（PBR）的多次镜面反射计算。最后，针对 1 SPP 极低采样率带来的蒙特卡洛高频噪声，完整集成了时空方差导向滤波（SVGF）降噪系统。通过几何运动矢量的时域历史重投影与动态方差引导的多轮 A-Trous 空域滤波，成功实现了高保真度的抗噪平滑处理。

性能测试与效果分析表明，本系统在 2K 原生分辨率下，能够以实时的性能开销（稳定大于 60 帧）输出平滑且物理正确的照片级渲染画面。本文的研究验证了以 Render Graph 为底座、结合光栅化、硬件光追与 SVGF 降噪的混合架构在现代底层图形引擎开发中的高效性与工程可行性。

\end{cnabstract}
\par
\vspace*{2em}

%%%%--  关键词 -----------------------------------------%%%%%%%%
\cnkeywords{Vulkan；实时光线追踪；混合渲染；Render Graph；基于物理的渲染 (PBR)；SVGF 降噪}

\newpage
%%%---- 3. 英文摘要 ------------------------------------%
\begin{enabstract}
\thispagestyle{empty}

In modern real-time computer graphics, traditional pure rasterization rendering has inherent physical limitations when dealing with global illumination, precise physical soft shadows, and complex specular reflections, while offline-level pure path tracing is difficult to meet the real-time frame rate requirements of interactive applications. To address this contradiction between "image quality and computing power", this paper designs and implements an industrial-grade high-performance real-time ray tracing hybrid rendering engine based on the modern explicit graphics API Vulkan 1.3.

The core work of this paper includes three aspects. First, a data-driven Render Graph scheduling architecture was designed to address the tedious GPU memory state synchronization pain points of Vulkan. Through topological analysis of a directed acyclic graph, the automatic derivation of pipeline barriers and optimization of memory reuse were achieved. Second, an efficient hybrid lighting pipeline was constructed to deeply decouple geometric rasterization from hardware ray tracing. Based on deferred rendering G-Buffer, combined with derivative face normal offset and Ray Query extension, precise and artifact-free soft and hard shadows were implemented. Furthermore, the RT Pipeline was used to comprehensively solve the Cook-Torrance microfacet model, completing multi-bounce specular reflection calculations based on physics (PBR). Finally, addressing the high-frequency Monte Carlo noise caused by the extremely low sampling rate of 1 SPP, a complete Spatiotemporal Variance-Guided Filtering (SVGF) denoising system was integrated. High-fidelity anti-noise smoothing was successfully achieved through temporal history reprojection of geometric motion vectors and multi-round spatial A-Trous filtering guided by dynamic variance.

Performance tests and effectiveness analysis show that this system can output smooth and physically correct photo-realistic rendering results with real-time performance overhead (stable above 60 FPS) at 2K native resolution. This research verifies the efficiency and engineering feasibility of a hybrid architecture based on Render Graph, combining rasterization, hardware ray tracing, and SVGF denoising in modern low-level graphics engine development.

\end{enabstract}
\par
\vspace*{2em}

\enkeywords{Vulkan; Real-time Ray Tracing; Hybrid Rendering; Render Graph; Physically Based Rendering (PBR); SVGF Denoising}

\makeatother % 恢复 @ 的原始定义